
\medskip
\phantomsection\label{par:sga3-VIA-5.5}%
\noindent\textbf{5.5.---}
Let \(G\) be a group locally of finite type and flat over an
Artinian local ring \(A\). We know from \SGARefLink{cor:sga3-VIB-2.3}{2.3} that the connected
component \(G^0\) of the identity is a normal open group subscheme
of \(G\), hence is also flat over \(A\). It then follows from \hyperref[thm:sga3-VIA-3.2]{3.2}
and \hyperref[par:sga3-VIA-5.2]{5.2} that
\[
  G^0\backslash G
\]
is an \(A\)-group scheme, flat over \(A\). Moreover, since every
connected component \(G^\alpha\) of \(G\) is saturated for the
equivalence relation defined by \(G^0\), the quotient
\(G^0\backslash G\) is the disjoint union of the
\(G^0\backslash G^\alpha\) (cf.\ \hyperref[par:sga3-VIA-4.3]{4.3}). In particular, the connected
component of the identity in \(G^0\backslash G\) is precisely
\[
  G^0\backslash G^0\simeq\operatorname{Spec}A,
\]
and hence the morphism
\[
  G^0\backslash G\longrightarrow\operatorname{Spec}A
\]
is a local isomorphism at the identity. Consequently,
\(G^0\backslash G\) is \'etale over \(\operatorname{Spec}A\), by
Expos\'e~VI\(_B\), \SGARefLink{prop:sga3-VIB-1.3}{1.3}.%
\footnote{\label{note:sga3-VIA-55}Editor's note: we have expanded
the preceding discussion from the original and added
\hyperref[prop:sga3-VIA-5.5.1]{Proposition~5.5.1} in order to make this result explicit.}
We therefore obtain the following proposition (for assertion~\SGARefLink{prop:sga3-VIA-5.5.1:item:ii}{(ii)},
see \textup{[DG70]}, \S~II.5, 1.7--1.10).

\medskip
\phantomsection\label{prop:sga3-VIA-5.5.1}%
\noindent\textbf{Proposition 5.5.1.---}
\emph{Let \(A\) be an Artinian local ring and \(G\) an \(A\)-group
locally of finite type and flat.
\begin{enumerate}
\item[\textup{(i)}] \(G^0\backslash G\) is an \'etale \(A\)-group.
\item[\textup{\SGAInlineRefTarget{prop:sga3-VIA-5.5.1:item:ii}(ii)}] Consequently, if \(A=k\) is an algebraically
closed field, \(G^0\backslash G\) is a constant \(k\)-group acting
simply transitively on the set of connected components of \(G\).
Thus, if \(G\) is algebraic, \(G^0\backslash G\) is finite.
\end{enumerate}}

\medskip
\phantomsection\label{par:sga3-VIA-5.6}%
\noindent\textbf{5.6.---}
Now let \(k\) be a perfect field and \(G\) a \(k\)-group locally of
finite type. We saw in \SGARefLink{sga3:VIA:0.2}{0.2} that \(G_{\mathrm{red}}\) is then a group
subscheme of \(G\). Moreover, the equivalence class of the identity
of \(G\) under the left action of \(G_{\mathrm{red}}\) on \(G\) is
the whole underlying space of \(G\). \hyperref[thm:sga3-VIA-3.2]{Theorem~3.2} therefore gives
the following result.

\medskip
\phantomsection\label{prop:sga3-VIA-5.6.1}%
\noindent\textbf{Proposition 5.6.1.---}
\emph{Let \(k\) be a perfect field and \(G\) a \(k\)-group locally
of finite type. Then the \(k\)-scheme
\[
  G_{\mathrm{red}}\backslash G
\]
is the spectrum of a finite local \(k\)-algebra with residue field
\(k\).}

\noindent
\footnote{\label{note:sga3-VIA-56}Editor's note: we have added the
number 5.6.1, as well as the proof that follows.}
Indeed, by \SGARefLink{thm:sga3-VIA-3.2}{3.2}, the quotient \(G_{\mathrm{red}}\backslash G\) has a
single point, with residue field \(k\), and is a \(k\)-scheme of
finite type. It is therefore the spectrum of a finite-dimensional
local \(k\)-algebra (cf.\ EGA~I, 6.4.4).

\medskip
\phantomsection\label{prop:sga3-VIA-5.6.2}%
\noindent\textbf{Proposition 5.6.2.---}
\emph{Let \(u:F\to G\) be a morphism between groups locally of
finite type over a perfect field \(k\). The following assertions
are equivalent:
\begin{enumerate}
\item[\textup{\SGAInlineRefTarget{prop:sga3-VIA-5.6.2:item:i}(i)}] \(u\) is flat.
\item[\textup{\SGAInlineRefTarget{prop:sga3-VIA-5.6.2:item:ii}(ii)}] The morphism
\[
  u^0:F^0\longrightarrow G^0
\]
is dominant, and the morphism
\[
  v:F_{\mathrm{red}}\backslash F
    \longrightarrow G_{\mathrm{red}}\backslash G
\]
induced by \(u\) is flat.
\end{enumerate}}

\noindent
\footnote{\label{note:sga3-VIA-57}Editor's note: we have expanded
the proof of \textup{\SGARefLink{prop:sga3-VIA-5.6.2:item:ii}{(ii)}} \(\Rightarrow\) \textup{\SGARefLink{prop:sga3-VIA-5.6.2:item:i}{(i)}} and
simplified the diagram below.}
Indeed, consider the following commutative diagram:
\begin{center}
  \phantomsection\label{diag:sga3-VIA-5.6.2-quotients}%
  \begin{tikzcd}[row sep=large,column sep=large]
    F
      \arrow[r,"p"]
      \arrow[d,"u"'] &
    F_{\mathrm{red}}\backslash F
      \arrow[d,"v"] \\
    G
      \arrow[r,"q"'] &
    G_{\mathrm{red}}\backslash G
  \end{tikzcd}
\end{center}
where \(p\) and \(q\) denote the canonical projections. By
\SGARefLink{concl:sga3-VIA-3.2-iv}{3.2(iv)}, \(p\) and \(q\) are faithfully flat. Consequently, if
\(u\) is flat, then
\[
  q\circ u=v\circ p
\]
is flat, and hence \(v\) is flat as well.

Conversely, suppose that \(v\) is flat and \(u^0\) is dominant.
Since \(u^0\) is quasi-compact (the scheme \(F^0\) is of finite type
over \(k\) by \SGARefLink{prop:sga3-VIA-2.4}{2.4}, and hence Noetherian), it maps the generic point
\(\xi\) of \(F^0\) to the generic point \(\eta\) of \(G^0\). Let
\(R\) be the finite local \(k\)-algebra such that
\[
  G_{\mathrm{red}}\backslash G=\operatorname{Spec}R,
\]
and let \(\mathfrak m\) be its maximal ideal. We have local
homomorphisms of local rings
\[
  R\longrightarrow\mathcal O_{G,\eta}
   \longrightarrow\mathcal O_{F,\xi}.
\]
Observe that the following square is Cartesian:
\begin{center}
  \phantomsection\label{diag:sga3-VIA-5.6.2-local-flatness}%
  \begin{tikzcd}[row sep=large,column sep=large]
    G_{\mathrm{red}}
      \arrow[r]
      \arrow[d] &
    G
      \arrow[d,"q"] \\
    \operatorname{Spec}(R/\mathfrak m)
      \arrow[r] &
    \operatorname{Spec}(R)
  \end{tikzcd}
\end{center}
It follows that
\[
  \mathcal O_{G,\eta}/\mathfrak m\mathcal O_{G,\eta}
    \simeq \mathcal O_{G_{\mathrm{red}},\eta}
    =\kappa(\eta),
\]
so that
\[
  \mathcal O_{F,\xi}/\mathfrak m\mathcal O_{F,\xi}
\]
is flat over
\(\mathcal O_{G,\eta}/\mathfrak m\mathcal O_{G,\eta}\).

On the other hand, since \(q\) and \(v\circ p\) are flat, \(G\)
and \(F\) are flat over \(R\). The local criterion for flatness
(cf.\ EGA~IV\(_3\), 11.3.10.2) therefore shows that
\(\mathcal O_{F,\xi}\) is flat over \(\mathcal O_{G,\eta}\); in
other words, \(u\) is flat at \(\xi\). It follows from \hyperref[lem:sga3-VIA-2.5.3]{2.5.3} that
\(u\) is flat.
